\todo[inline]{combine with Quack. show the difference between specialized oracles and master-problems vs just logit and gurobi}
\todo[inline]{add results for all games}

In this section we present an algorithm that is much more widely applicable at the cost of relaxing the requirement on the exactness of solutions from the previous section. We focus on the computation and approximation of mixed strategy equilibria for the whole class of multiplayer general-sum continuous games. To solve such generic games, we vastly extend the scope of applicability of the double oracle algorithm, initially designed and proved to converge only for two-player zero-sum games. Specifically, we propose an iterative strategy generation technique, which splits the original problem into the master problem with only a finite subset of strategies being considered, and the subproblem in which an oracle finds the best response of each player. This simple method is guaranteed to recover an approximate equilibrium in finitely many iterations.

Further, we argue that the Wasserstein distance (the earth mover's distance) is the right metric for the space of mixed strategies for our purposes. Our main result is the convergence of this algorithm in the Wasserstein distance to an equilibrium of the original continuous game. The numerical experiments show the performance of our method on several classes of games including randomly generated examples.

The double oracle algorithm \cite{McMahan03} was extended from finite games and proved to converge for all two-player zero-sum continuous games \cite{Adam21}. The algorithm is relatively straightforward. It is based on the iterative solution of finite subgames and the subsequent extension of the current strategy sets with best response strategies. Despite its simplicity, this method was successfully adapted to large extensive-form games \cite{Bosansky14}, Bayesian games \cite{Li21}, and security domains with complex policy spaces \cite{Xu21}.

In this contribution, we extend the double oracle algorithm beyond two-player zero-sum continuous games \cite{Adam21}. This is a significant extension since we had to design a new convergence criterion to track the behavior of the algorithm. On top of that, there are relatively few methods for computing or approximating equilibria of continuous games, which narrows down the continuous games for which the exact solution can be found and used as a benchmark in our experiments. Our main result guarantees that the new method converges in the Wasserstein distance to an equilibrium for any general-sum $n$-player continuous game. Interestingly enough, it turns out that the Wasserstein distance represents a very natural metric on the space of mixed strategies. We demonstrate the computational performance of our method on selected examples of games appearing in the literature and on randomly generated games.